\section{The shared graph comparison and reflected target grids}\label{sec:graph}

This section proves the common graph argument for
Theorem~\ref{thm:model-comparison}, then the residue-energy estimate
needed for its reflected family. The fixed-shift major arcs are treated
separately in Section~\ref{sec:fixed-shift-comparison}. No arithmetic
form of the indicators is used here beyond \textup{(C1)--(C3)}.
Two structural points matter. The labels lie in a fixed reduced class,
so the character frame must identify restrictions differing by a known
unit phase. A target-sized displacement requires retaining its exact
rational grids in the energy estimate.

Write \(L=\log x\).  There are \(\asymp\log L\) small banks and big
banks, with primes respectively in
\[
 [\exp(L^{1/8}),\exp(L^{1/4})],\qquad
 [\exp(L^{1/3}),\exp(L^{2/5})].
\]
Every label is \(c\bmod L_0\); primes dividing the nonzero displacement
\(k\) are removed. We use \(k=T=x\) for sums and fixed \(k=h\)
for shifts. All banks are disjoint. Their laws are
\(\mu_g(p)=(V_gp)^{-1}\), where \(V_g=\sum_{p\in\mathcal P_g}p^{-1}\)
is bounded above and below by positive fixed constants.
Banks in this section may be any active subfamilies of the banks
constructed from the class in \textup{(C2)}, provided that each size has
at least a fixed positive multiple of \(\log L\) banks.
All constants may depend on these fixed bounds and on \(L_0\).

\subsection{The character frame for progression banks}

The ideal-kernel construction in \cite[Section~4]{weightedgraphs}
uses primitive characters of conductor at most \(Q=L^{q_*}\).
On labels \(c\bmod L_0\), this list is redundant up to phases.
For example, the principal character and \(\chi_{-4}\) agree on
primes \(1\bmod4\) and are opposite on primes \(3\bmod4\).
The quotient below removes this redundancy while retaining its phase.

Let \(\mathcal X_Q\) be the primitive characters of conductor at most
\(Q\), and define
\[
 \eta\sim\eta'
 \quad\Longleftrightarrow\quad
 \operatorname{cond}(\eta\overline{\eta'})\mid L_0.
 \tag{5.1}\label{eq:gr-equivalence}
\]
The conductor here belongs to the primitive character inducing the
quotient on a common modulus.  Choose one representative from each
class and denote the resulting set by \(\mathcal R_Q\).
Thus \(|\mathcal R_Q|\le Q^2\).

To verify the equivalence and its meaning, put
\(M_0=\lcm(L_0,\operatorname{cond}\eta,\operatorname{cond}\eta')\).
Reduction from \((\Z/M_0\Z)^\times\) onto
\((\Z/L_0\Z)^\times\) is surjective.  Its kernel consists of the
units \(1\bmod L_0\).  A character is trivial on that kernel exactly
when it factors through the latter group.  This is equivalent to
its primitive conductor dividing \(L_0\). Consequently, if
\(\eta\sim\eta'\) and their quotient is induced by
\(\vartheta\bmod L_0\), then on a product of \(r\) labels,
\[
                 \eta(D)=\vartheta(c)^r\eta'(D).
\]
Both comparison products have the same slot type and the same phase.
Inequivalent characters have zero average on the kernel, and also
on any fixed coset of it: translation multiplies their sum by a unit phase.

Fix a nonempty set \(\Lambda\) of at most \(m\) slots per big bank,
where \(m\) is fixed before the eventual graph mark length.
Let \(D_\Lambda\) be the product of independent labels in these slots.
Partition its logarithm into cells
\(I_j=[jh,(j+1)h)\), and write
\(\nu_j=\Prob(\log D_\Lambda\in I_j)\).
For inequivalent \(\eta,\eta'\in\mathcal R_Q\), and any prescribed
\(A>0\), Siegel--Walfisz gives
\[
 \left|\E\big[
  \1_{\log D_\Lambda\in I_j}
        \eta(D_\Lambda)\overline{\eta'(D_\Lambda)}\big]\right|
       \ll_A L^{-A}.                                  \tag{5.2}\label{eq:gr-cell}
\]
Here and below the accuracy in Siegel--Walfisz is chosen after all
fixed logarithmic exponents.

Indeed, freeze all but one label.  Its allowable values form an
interval inside its original bank interval, still restricted to
\(c\bmod L_0\). On a common modulus at most \(L_0Q^2\), the sum of
\(\eta\overline{\eta'}\) over the appropriate reduced residue classes
is zero.  Siegel--Walfisz, summed over these classes, therefore gives
\[
 \sum_{\substack{p\le z\\p\equiv c\ (L_0)}}
       \eta(p)\overline{\eta'(p)}
       \ll_K L_0Q^2z(\log z)^{-K}.
\]
Partial summation bounds the reciprocal-prime sum on any subinterval
of \([\exp(L^{1/3}),\exp(L^{2/5})]\) by
\(O_K(L_0Q^2(L^{1/3})^{1-K})\).
This is an absolute bound; no lower bound for the subinterval's
length or harmonic mass is needed.  Divide by \(V_g\), then average
the frozen labels.  Deleting target divisors changes the answer by at most
\[
 \sum_{\substack{p\mid k\\p\ge\exp(L^{1/3})}}\frac1p
 \le \frac{\log(2|k|)}{L^{1/3}}\exp(-L^{1/3}),               \tag{5.3}\label{eq:gr-deletion}
\]
which is smaller than every fixed inverse power of \(L\).
This proves \eqref{eq:gr-cell} uniformly for the actual banks.

On cells with \(\nu_j\ge\xi\), define
\[
 K_\Lambda(D,X)=
 \sum_{\nu_j\ge\xi}
 \frac{\1_{\log D\in I_j}\1_{\log X\in I_j}}{\nu_j}
 \sum_{\eta\in\mathcal R_Q}\overline{\eta(D)}\eta(X).
 \tag{5.4}\label{eq:gr-kernel}
\]
Both product variables have the same law.  The elementary bounds are
\[
 |K_\Lambda|\le Q^2\xi^{-1},\qquad
 \sup_X\E_D|K_\Lambda(D,X)|\le Q^2,\qquad
 \#\{j:\nu_j>0\}\ll_m1+(\log L)L^{2/5}/h.              \tag{5.5}\label{eq:gr-kernel-bounds}
\]
The last inequality uses only the range of the logarithm.

On a retained cell, the kernel is the positive frame operator
\[
 f\longmapsto
 \sum_{\eta\in\mathcal R_Q}\eta\,\langle\eta,f\rangle
\]
for the conditional product probability measure.
Its Gram matrix has diagonal entries one, since all labels exceed
\(L_0Q^2\), and off-diagonal entries
\(O_A(\xi^{-1}L^{-A})\) by \eqref{eq:gr-cell}.
Thus its norm is at most
\(1+O_A(Q^2\xi^{-1}L^{-A})\le2\).
Distinct cells are orthogonal blocks, so
\[
             \|K_\Lambda\|_{2\to2}\le2.                \tag{5.6}\label{eq:gr-frame-norm}
\]
This also holds for tuple tests: conditional expectation from tuples
to their product is a contraction.
If \(\psi\) is a character of modulus \(r\le Q\), its primitive
inducing character is equivalent to one member of \(\mathcal R_Q\).
For \(r_\Lambda\) selected slots and \(\psi=\vartheta\eta\) on
these units, that summand gives
\(\vartheta(c)^{r_\Lambda}\eta(X)=\psi(X)\). The other summands give
\[
 \E_D K_\Lambda(D,X)\psi(D)
       =\psi(X)+O_A(Q^2\xi^{-1}L^{-A})                 \tag{5.7}\label{eq:gr-reproduce}
\]
at every retained \(X\).  This is where quotienting the frame is
essential: the unquotiented frame can reproduce the same constant
several times.

\begin{lemma}[Progression ideal comparison]\label{lem:gr-ideal}
For every prescribed ideal accuracy \(E_{\rm id}\) and fixed phase
range \(|\zeta|\le L^{C_4}\), the progression banks admit the signed
comparison family of \cite[Theorem~4.1]{weightedgraphs}.
Its ideal norm is at most \(L^{-E_{\rm id}}\), uniformly in every
real additive frequency.  Its coefficients ignore small-bank labels
and shared labels, free at most \(m\) slots per big bank, and have
total supremum at most \(L^{C_1}\).
The parameters \(m,C_1\) and all kernel powers are fixed before the
eventual mark length.
\end{lemma}

\begin{proof}
We give the comparison estimates used in the construction, including
the only step where restricting the prime labels changes the proof.

For a product formed from \(k_g\le m\) draws in each of at most
\(C_0\log L\) disjoint banks, unique factorization gives its exact law
\[
 n\Prob(D_\Lambda=n)
   =\prod_g\frac{k_g!}{V_g^{k_g}\prod_{p\in\mathcal P_g}a_p!},
 \qquad
 \Prob(D_\Lambda=n)\le L^{K_0}/n,                     \tag{5.8}\label{eq:gr-productlaw}
\]
where
\(K_0=C_0(\log(m!)+m\log\max(1,v_-^{-1}))\) and
\(V_g\ge v_->0\).
The factor \(C_0\) multiplies both terms of this exponent.
For a tuple test, Cauchy--Schwarz on each product fibre gives an
integer coefficient sequence of squared norm at most
\(L^{K_0}/U\) times the tuple squared norm on a product dyad
\([U,2U)\).

The elementary integer bilinear estimate
\cite[Lemma~4.2]{weightedgraphs}, or \cite[Lemma~6.2]{walsh}, gives
\[
 \left|\sum_{u,v}a_ub_v\ee{\theta uv}\right|^2
 \ll \|a\|_2^2\|b\|_2^2
             (1+U/r)(V+r\log(2r))                  \tag{5.9}\label{eq:gr-integer-bilinear}
\]
when the two supports lie in containing intervals of lengths
\(O(U),O(V)\), \((a,r)=1\), and
\(|\theta-a/r|\le r^{-2}\).
For completeness, Cauchy--Schwarz and expansion leave
\(V+\sum_{1\le j\ll U}\min(V,\|\theta j\|^{-1})\).
Divide the \(j\)'s into blocks of diameter less than \(r/2\).
Within a block the rational parts are \(1/r\) separated and their
perturbations differ by at most \(1/(2r)\); hence the sum is
\(O(V+r\log(2r))\) per block.  This proves
\eqref{eq:gr-integer-bilinear} and permits arbitrary subsets of the
containing intervals.

Apply this estimate to two nonempty selected product sets \(A,B\).
On product dyads \(U,V\), use Dirichlet approximation with denominator
ceiling \(\lfloor UV/S\rfloor\).  When the resulting denominator
exceeds \(Q\), \eqref{eq:gr-productlaw} gives the normalized bound
\[
 \ll L^{K_0}\big(o(L^{-A})+Q^{-1}+L^{7/5}/S\big)^{1/2}.
 \tag{5.10}\label{eq:gr-highden}
\]
There are at most \(L^{7/5}\) product dyads per variable.
The same bound applies to the kernel terms by moving each kernel
onto its endpoint test using \eqref{eq:gr-frame-norm}.
The factor \((uv)^{i\zeta}\) separates and preserves test norms.

For denominators \(r\le Q\), all products are units modulo \(r\).
The multiplicative-character expansion of the rational phase has
total absolute coefficient at most \(\sqrt{\varphi(r)}\).
Replace each primitive character by its representative in
\(\mathcal R_Q\) and apply \eqref{eq:gr-reproduce}.
Grouping characters after inserting their unit phases cannot increase
coefficient mass, so the reproduction error is
\(O_A(Q^{5/2}\xi^{-1}L^{-A})\).
Write the remaining phase as
\(M(v)=\exp(i\zeta v)\ee{\beta\exp v}\).
On matched cells its derivative is
\(O(L^{C_4}+S)\), and the two logarithmic products move by at most
\(2h\).  By \eqref{eq:gr-kernel-bounds}, this costs
\(O(Q^4h(S+L^{C_4}))\).
The union of discarded cells has probability at most
\((N_A+N_B)\xi\), where \(N_A,N_B\) count the possible cells.
Cauchy--Schwarz bounds its contribution by the square root of this
probability times the two test norms.
These estimates prove the two-product comparison, for arbitrary
\(L^2\) tuple tests, to any prescribed power \(L^{-A_0}\).

There is no circular choice of precisions.  One possible hierarchy is
\[
 \begin{split}
 B_0&>A_0+2(7/5)+10,\\
 q_*,s_*&>2(B_0+K_0+10)+7/5,\\
 h_*&>4q_*+\max(s_*,C_4)+B_0+10,\\
 x_*&>2B_0+7/5+h_*+10,
 \end{split}
\]
with \(Q=L^{q_*},S=L^{s_*},h=L^{-h_*},\xi=L^{-x_*}\).
Choose the accuracy in \eqref{eq:gr-cell} afterwards.
All these choices are independent of the eventual mark length.

Here is the signed cancellation which completes the ideal construction.
Split the big banks into two blocks of comparable size and designate
\(m\) probe slots per group.  For nonempty probe subsets
\(A,B\) in the two blocks, free their slots and give the comparison
coefficient
\[
              -(-1)^{|A|+|B|}K_A(D_A,X_A)K_B(D_B,X_B).
 \tag{5.11}\label{eq:gr-signed}
\]
Here \(D_A,D_B\) are auxiliary-label products and \(X_A,X_B\)
are the unshared endpoint products.  The coefficient has no small
or shared labels.  There are at most \(2^{mC_0\log L}\) comparisons;
by \eqref{eq:gr-kernel-bounds} their total supremum is at most
\[
        1+2^{mC_0\log L}Q^4\xi^{-2}\le L^{C_1}.         \tag{5.12}\label{eq:gr-coeffcost}
\]
Expanding both kernels also separates each coefficient into a sum
of products of four bounded factors: one on each auxiliary product
and one on each unshared endpoint list.  The number of terms and
total scalar coefficient mass are fixed powers of \(L\).

We now make the exceptional projection and the cancellation explicit.
Write the two big-bank blocks as \(\mathcal B_1,\mathcal B_2\), of
sizes \(s_1,s_2\), and their designated probe coordinates as
\(\mathcal I_1,\mathcal I_2\).  For a coordinate \(i\), let \(E_i\)
average its independent label.  These commuting operators are
orthogonal projections.  Let \(Q_{\rm tail}\) average the last
\(\ell\) input or target coordinates in every group, and let
\(\Pi\) be within-group symmetrization on all \(M=J+\ell\) coordinates.
The ideal multipliers do not use those last coordinates.
For symmetric endpoint tests \(F,H\), put
\(F'=Q_{\rm tail}F,\ H'=Q_{\rm tail}H\).
Both averaging operations are contractions.

Let \(D_{\rm sh}\) be the product of the first \(J\) coordinates,
and write
\[
 \Phi_{\Theta,\zeta}
       =D_{\rm sh}^{i\zeta}\ee{\Theta D_{\rm sh}},\qquad
 R(F',H')=\E\big[F'\overline{H'}\Phi_{\Theta,\zeta}\big].
\]
Thus \(R\) is a bilinear pairing of norm at most one.
Resolve the probes in the first block of \(F'\) and the second
block of \(H'\):
\[
 \begin{split}
 F'&=\sum_{S\subset\mathcal I_1}F_S,\qquad
 F_S=\prod_{i\in S}(I-E_i)
             \prod_{i\in\mathcal I_1\setminus S}E_iF',\\
 H'&=\sum_{T\subset\mathcal I_2}H_T,\qquad
 H_T=\prod_{i\in T}(I-E_i)
             \prod_{i\in\mathcal I_2\setminus T}E_iH'.
 \end{split}
\]
These are orthogonal decompositions.  Define
\(P_1=\prod_{i\in\mathcal I_1}(I-E_i)\),
\(P_2=\prod_{i\in\mathcal I_2}(I-E_i)\), so that
\(F_{\rm full}=P_1F'\), \(H_{\rm full}=P_2H'\).

In one group the original symmetric vector decomposes orthogonally
over subsets of its \(M\) coordinates.  Symmetry makes the squared
norms equal among subsets of cardinality \(j\).
Requiring all \(m\) probes to occur and all last \(\ell\) slots
to be absent therefore retains exactly the fraction
\[
 \frac{\binom{M-m-\ell}{j-m}}{\binom Mj}
     =\frac{(j)_m(M-j)_\ell}{(M)_{m+\ell}}
     \le 2(j/M)^m(1-j/M)^\ell
     \le2\ell^\ell m^{-\ell}=:\theta_m
\]
when \(M\ge(m+\ell)(m+\ell-1)\); impossible cardinalities give zero.
For a vector which is not a tensor product across groups, first
resolve by the tuple of cardinalities in the groups of a block.
Independent within-group permutations make all subset tuples of
those cardinalities have equal squared norms.  The retained fraction
is consequently the product of the displayed fractions.
Summing over the cardinality tuples proves the operator estimates
\begin{equation}\label{eq:gr-blockprojection}
 \|P_1Q_{\rm tail}\Pi\|\le\theta_m^{s_1/2},\qquad
 \|P_2Q_{\rm tail}\Pi\|\le\theta_m^{s_2/2}.
\end{equation}
Averaging tails outside the block is a further contraction.
We use \(\ell=1\), and each \(s_i\ge c_0\log L/3\).

For a component pair \(F_S,H_T\), put
\(U=\mathcal I_1\setminus S\), \(V=\mathcal I_2\setminus T\).
A comparison freeing \(A\subset\mathcal I_1\),
\(B\subset\mathcal I_2\) vanishes unless \(A\subset U,B\subset V\).
Indeed, if a freed input coordinate belongs to \(S\), it is
mean zero, absent from the target, and absent from the comparison
multiplier: integrating it gives zero.  The corresponding statement
holds for a freed target coordinate belonging to \(T\).
This uses the orientation in \eqref{eq:gr-signed}: the \(A\)
endpoint labels are from the target and the \(B\) endpoint labels
are from the input.

If \(A\subset U,B\subset V\), condition on all shared coordinates
outside \(A\cup B\), with product \(C\).
The input test is independent of the freed \(A\) labels and the
target test is independent of the freed \(B\) labels.
The two-product kernel comparison established above therefore
replaces the unsigned comparison by \(R(F_S,H_T)\), with error
\[
                    L^{-A_0}\|F_S\|_2\|H_T\|_2.
\]
Its frequency is \(\Theta C\); uniformity for every real frequency
is essential.  Cauchy--Schwarz over the conditioned coordinates
gives exactly the displayed unconditional norms.

For nonempty \(U,V\), the total unsigned subset sign is
\[
 \sum_{\substack{\varnothing\ne A\subset U\\
                 \varnothing\ne B\subset V}}
       (-1)^{|A|+|B|}
 =\left(\sum_{\varnothing\ne A\subset U}(-1)^{|A|}\right)
  \left(\sum_{\varnothing\ne B\subset V}(-1)^{|B|}\right)=1.
\]
The initial minus sign in \eqref{eq:gr-signed} cancels the raw
pairing.  If either missing set is empty, no comparison survives.
All exceptional raw component pairs together have the exact identity
\[
 R_{\rm exc}(F,H)=
 R(P_1F',H')+R((I-P_1)F',P_2H').
\]
In operator notation, if \(M_\Phi\) multiplies by
\(\Phi_{\Theta,\zeta}\), their operator on the symmetric space is
\[
 T_{\rm exc}=
 \Pi Q_{\rm tail}
   \big[M_\Phi P_1+P_2M_\Phi(I-P_1)\big]
                 Q_{\rm tail}\Pi .
\]
Consequently
\[
 \|T_{\rm exc}\|
       \le\theta_m^{s_1/2}+\theta_m^{s_2/2}.
\]
This identity sums the exceptional components before taking an
absolute value.  It is the reason their bound has no subset-count
or comparison-coefficient factor.

There are at most \(3^{P}\) admissible component/free-subset triples,
where \(P=m(s_1+s_2)\le mC_0\log L\): each probe is indicated,
missing but not freed, or freed.  Thus the full signed ideal
operator satisfies
\[
 \|T_{\rm signed}-T_{\rm exc}\|
           \le3^{P}L^{-A_0}.
\]
Choose \(m\) so that \(\theta_m<1\) and
\(-c_0\log\theta_m/6>E_{\rm id}+3\), and then choose
\(A_0>E_{\rm id}+mC_0\log3+10\).
The kernel precisions above and the exponent \(C_1\) in
\eqref{eq:gr-coeffcost} are chosen after these choices.
Only afterwards require \(J\ge m\) and
\(J+\ell\ge(m+\ell)(m+\ell-1)\); increasing this fixed \(J\)
to meet the physical root-moment theorem changes no kernel
parameter or coefficient exponent.
The resulting ideal norm is at most \(L^{-E_{\rm id}}\).
This proves the asserted ideal norm and coefficient conditions.
\end{proof}

\subsection{A large displacement in the physical graph}

\begin{lemma}[Large-target root moment and pairing]\label{lem:gr-large}
The root moment and endpoint pairing from
\cite[Theorem~3.5 and Corollary~3.11]{weightedgraphs} hold for any
nonzero integer \(k\) coprime to the bank primes, with the same
logarithmic accuracies.  In particular they hold for \(k=T\asymp x\)
after target-prime deletion.  The first endpoint is required to be
independent of its marks, exactly as in the imported pairing.
\end{lemma}

\begin{proof}
Fix the labels, patterns, and permutations of a closed walk of length
\(2R\), where \(R=\lceil L^{0.52}\rceil\).
Its position offsets are \(ks_i\).  Return to the initial integer
is equivalent to \(s_{2R}=0\), independent of the nonzero \(k\).
The root dependence consists of divisibility and damping conditions
at the offsets \(ks_i\).

The independent-residue replacement
\cite[Lemma~3.6]{weightedgraphs} is uniform in the numerical values
of these offsets.  Indeed its active modulus has logarithm
\(O_M(R(\log L)L^{2/5})=o(L)\).
The Bonferroni truncation has the same modulus and term count,
and each interval CRT error is \(O(1)\), independently of the
specified residues.  Once the deterministic damping from active
primes is factored out, the error is relative and can be summed with
absolute path coefficients using the imported absolute row bound.

In the resulting independent residue space, replace the root residue
\(z_p\) by \(k^{-1}z_p\bmod p\) for every bank prime \(p\).
Every fixed-path divisibility and damping expectation is then exactly
the expectation for \(k=1\).
The remaining signed coefficients are unchanged.
Thus the \(k=1\) root moment gives the identical bound for this \(k\).
Negative root intervals follow either from the same CRT count or
from \(n\mapsto-n,\ k\mapsto-k\).

It remains to avoid a large factor in the endpoint pairing.
Let \(D_{\max}\) bound the edge products, so that
\(\log D_{\max}=O_M((\log L)L^{2/5})\).
Partition the first-endpoint interval \(I\) into blocks of length
\(\lceil |k|D_{\max}\rceil\), as in the source pairing, and refine
each block by residue modulo \(|k|\).
The operator preserves these residues.
Each refined block has \(O(D_{\max}+1)\) integer positions.
For \(B=\mathcal A_k\mathcal A_k^*\), the matrix
\((\langle u_n,B^Ru_m\rangle)\) on one such block is positive
semidefinite.  Its largest eigenvalue is at most its trace.
Here \(u_n\) is the unnormalized vector equal to one on all marks at
\(n\).  Mark independence of the first endpoint \(f_n\) gives
\[
 \langle f_j,B^Rf_j\rangle
 \le \Big(\sum_{n\in I_j}|f_n|^2\Big)
             \sum_{n\in I_j}\langle u_n,B^Ru_n\rangle.
\]
Spectral H\"older and Cauchy--Schwarz give the source block estimate;
its only dimension factor is now
\[
 \big(O(D_{\max})\sup|f_n|^2\big)^{1/(2R)}
 \le
 \exp\!\left(O_M\!\left(
 \frac{(\log L)L^{2/5}+\sqrt L}{L^{0.52}}\right)\right)=O(1).
 \tag{5.13}\label{eq:gr-block}
\]
Second-endpoint neighborhoods have bounded overlap within a residue
class and are disjoint between residue classes.
H\"older over blocks uses only the sum of root traces on the whole
original interval, bounded by the root theorem.
No CRT estimate on an individual residue class modulo \(k\) is used.
This proves the pairing uniformly in \(k\).
\end{proof}

For bounded invariant endpoint functions, the list moments needed
here do not create a mark-dependent logarithmic loss.  For a fixed
ordered distinct list with product \(P\), divisibility gives
\(O(X/P)\) integers in an interval of length \(O(X)\) contained in
\([-CX,CX]\).  For fixed mark length, \(P=x^{o(1)}\).
After summing the list weights,
\[
 \sum_{\text{lists}}\frac1P\prod_gV_g^{-M}\le1.
 \tag{5.14}\label{eq:gr-list-moments}
\]
Thus the endpoint \(r\)-moments are \(O_r(X)\).
Condition \textup{(C1)} bounds every fixed-depth model without a
depth-dependent logarithmic exponent. Thus this observation is valid
before the final depth is chosen.

\subsection{Shared labels and the two comparison geometries}

For sums take \(k=x\); for shifts take fixed \(k=h\ne0\).
We first obtain the same comparison identity in both cases.

We spell out the arithmetic consequence of the pairing, since
the large reflected target is not covered by the fixed-shift theorem
of \cite{walsh}.
For a vector \(\mathbf t=(t_g)\) and a bounded function \(b\) on its
ordered distinct lists, put
\[
 W_{\mathbf t}^b(n)=
 \prod_g\big(\1_{\omega_g(n)=t_g}V_g^{-t_g}\big)
       \sum_{\substack{\mathbf p\text{ a }\mathbf t\text{-list}\\
                        p_{g,j}\mid n}} b(\mathbf p).
 \tag{5.15}\label{eq:gr-listweight}
\]
Here \(\omega_g(n)\) counts distinct bank-prime divisors of \(|n|\).
For fixed \(t_g\le m+1\), these weights are bounded by a fixed
power of \(L\).  In particular \(W_{\mathbf1}^1=\Omega_{\mathcal P}\).

Let \(\alpha,\beta\) be bounded functions invariant under multiplying
by every bank prime. For sums extend the first centered function to
the negative half-line by absolute value and the second function to
the positive half-line. For shifts use positive endpoints on both
sides. These sign restrictions preserve prime-multiplication invariance.
Take a smooth \(\Psi\) supported in a compact subset of \((0,\infty)\).
The pairing and Lemma~\ref{lem:gr-ideal} imply
\[
 C_0+\sum_\nu c_\nu C_\nu=O_A(L^{-A}),\qquad
 \#\{\nu\}+\sum_\nu|c_\nu|\le L^{C},                  \tag{5.16}\label{eq:gr-comparison}
\]
where
\[
 C_0=\sum_{w\ne0}\frac{\Psi(|w|/x)}{|w|}
       \alpha(w)\beta(w+k)
       W_{\mathbf1}^1(w)W_{\mathbf1}^1(w+k)
\]
and each comparison has the form
\[
 \begin{split}
 C_\nu=\E_{D_1,D_2}a_\nu(D_1)b_\nu(D_2)
 \sum_{w\ne0}\frac{\Psi(|w|/(xD_1D_2))}{|w|}
  \alpha(w)\beta(w+kD_1D_2)\\
 {}\times W_{\mathbf t}^{u_\nu}(w)
                  W_{\mathbf t}^{v_\nu}(w+kD_1D_2).
 \end{split}                                                     \tag{5.17}\label{eq:gr-comparison-term}
\]
The two auxiliary products are independent, use at most \(m\) draws
per big bank, and each contains at least one big-bank prime.
The four separated coefficient functions have modulus at most one.
The list coefficients \(u_\nu,v_\nu\) ignore all small-bank labels;
in every small bank \(t_g=1\).

Here is the derivation of these precise assertions from the physical
graph.  A pattern has \(M-1\) dilation slots per bank, some shared
and some freed, and one distinguished unshared endpoint slot.
Write its edge product as \(D=D_RD_C\), with shared product \(D_R\)
and auxiliary product \(D_C\).
For a source of either sign in \(|n|\in[Z,2Z)\), an originally nonzero edge
weight satisfies
\(\Psi(|n|/(xD))\ne0\), hence \(D\ll Z/x\).
Since \(|k|\ll x\), this implies \(|n+kD|\ll Z\).
Insert a second-endpoint cutoff on \([-C_\Psi Z,C_\Psi Z]\)
before Mellin inversion.  It equals one on every contributing edge.
After inversion retain this same cutoff.  An edge originally having
zero \(\Psi\) still has zero full Mellin integral, so the insertion
does not change the identity.  No cutoff involving shared-label
products has been introduced.

Mellin inversion of \(\Psi\), with the harmonic weight \(1/|n|\),
now applies the pairing of Lemma~\ref{lem:gr-large} on each of
\(O(L)\) source dyads.  Its tails are bounded by the absolute
row/column estimate, its main part by the small ideal norm, and
its endpoint norms by \eqref{eq:gr-list-moments}.
Compressing the physical operator to exact-count vertices
\(\omega_g(n)=\omega_g(n')=M\) is legitimate:
the corresponding projection depends only on the integer, preserves
mark independence, and sets both damping factors equal to one.

For completeness, put \(n=D_Rw\), \(n'=D_Rw'\), so
\(w'=w+kD_C\), and put \(t_g=1+\#\{\text{freed slots in }g\}\).
Invariance and harmonic scaling give
\[
 \alpha(D_Rw)=\alpha(w),\quad \beta(D_Rw')=\beta(w'),\quad
 \frac{\Psi(|n|/(xD))}{|n|}
       =\frac1{D_R}\frac{\Psi(|w|/(xD_C))}{|w|}.
\]
The graph normalization at a shared prime combines with \(1/D_R\)
to give its probability law \(\mu_g\).
More explicitly, with \((D_R)_g\) its product in bank \(g\),
\[
 \frac{V_g^{-M-t_g}}{(D_R)_g}
       =V_g^{-2t_g}
           \prod_{\text{shared slots in }g}\mu_g(p).   \tag{5.18}\label{eq:gr-sharednormal}
\]
For fixed reduced endpoints and fixed auxiliary labels,
the shared-label probability is
\[
 \prod_g\1_{\omega_g(w)=\omega_g(w')=t_g}
              +O(L^{C_M}e^{-L^{1/8}}).               \tag{5.19}\label{eq:gr-sharedprob}
\]
Indeed there are only \(O_M(\log L)\) draws, each atom is
\(O(e^{-L^{1/8}})\), and \(ww'D_C\) has \(O(L)\) prime divisors.
If the reduced counts are \(t_g\), distinct draws avoiding those
divisors reconstruct the exact graph edge.  If a reduced count is
larger, a valid reconstruction must choose a shared prime already
dividing that endpoint, an event of the displayed small probability.
This proves \eqref{eq:gr-sharedprob}, also for negative \(w\).

Summing its absolute error costs only fixed logarithmic powers,
by bounded endpoints and the harmonic weight.
The restrictions excluding auxiliary labels from endpoint lists
can then be removed: for an auxiliary prime \(p\mid D_C\),
\(p\mid w\) if and only if \(p\mid w+kD_C\); these positions have
density \(O(1/p)\) in the containing interval.
Every auxiliary prime is at least \(\exp(L^{1/3})\).
An unshared prime common to the two endpoints divides \(kD_C\);
bank-prime deletion implies it divides \(D_C\), so the same estimate
covers it.  Thus these errors remain smaller than every logarithmic
power.  Finally expand the separated coefficients from
\eqref{eq:gr-signed}.  Formula \eqref{eq:gr-sharednormal} yields
exactly \eqref{eq:gr-listweight}, while \(D_C=D_1D_2\).
The raw pattern has \(D_C=1,\mathbf t=\mathbf1\).
This proves \eqref{eq:gr-comparison}--\eqref{eq:gr-comparison-term}
with their stated independence properties.

\subsection{Prime sums at target-sized conductor}

Only this subsection uses the zero-free input.
It is important that its conductor range is polynomial in \(x\),
rather than a power of \(\log x\).

\begin{lemma}[Prime factor estimate]\label{lem:gr-primes}
Assume the fixed all-conductor zero-free half-plane stated in
Section~\ref{sec:inputs}.
Fix \(\sigma_1\in(11/12,1)\).
For \(q\le x^B\), every character \(\chi\bmod q\), and
\(P\ge\exp(L^a)\), uniformly on subintervals \(I\subset[P,2P]\),
\[
 \sum_{p\in I}\chi(p)p^{it}
 =\delta_\chi\int_I\frac{u^{it}}{\log u}\,du
       +O((1+|t|)P^{\sigma_1}L^C).                  \tag{5.20}\label{eq:gr-primes}
\]
Here \(\delta_\chi=1\) for a principal character and zero otherwise,
and the estimate holds for every real \(t\).
For progression banks the main vanishes unless \(\chi\) is induced
from a character modulo \(L_0\).
Removing target or modulus divisors preserves the error bound.
\end{lemma}

\begin{proof}
A fixed distance to the right of the zero-free boundary, the local
logarithmic-derivative formula and functional equation give
\[
 \left|\frac{L'}{L}(s,\chi)+\frac{\delta_\chi}{s-1}\right|
                 \ll \log(q(|\Im s|+3)).
 \tag{5.21}\label{eq:gr-logderivative}
\]
Imprimitive Euler factors contribute \(O(\log q)\), since the real
part is fixed and positive.
Apply truncated Perron to
\(-L'/L(s,\chi)\), at
\(c=1+1/\log Z\), height \(Z^2\), and a half-integer endpoint
\(Z\in[P/2,4P]\).
The truncation error is \(O(\log^2 Z)\): near \(Z\), sum the bounds
\(O(Z\log Z/(Z^2|n-Z|))\); away from \(Z\), use absolute convergence
at \(c\).  Shift to \(\sigma_1\).
The principal pole gives \(\delta_\chi Z\), the new vertical integral
is \(O(Z^{\sigma_1}\log^2(qZ))\), and horizontal edges cost
\(O(Z^c\log(qZ)/(Z^2\log Z))\).
Prime powers cost \(O(Z^{1/2}\log^2 Z)\).
Replacing a real endpoint by the adjacent half-integer changes its
prime-weighted count by \(O(\log Z)\).  Therefore
\[
       \sum_{p\le Z}\chi(p)\log p
          =\delta_\chi Z+O(Z^{\sigma_1}L^C).
\]
Partial summation against \(u^{it}/\log u\) proves
\eqref{eq:gr-primes} with its linear \(1+|t|\) growth.

Expand the progression condition into the finitely many characters
modulo \(L_0\).  A constituent is principal exactly when the original
character is induced from \(L_0\).
There are at most \(O(\log(q(2+|k|)))=O(L)\) deleted prime divisors;
their absolute unweighted contribution is absorbed by the error.
\end{proof}

In particular, for a nonexceptional character the prime factor saves
\(P^{\sigma_1-1}\) against its trivial length.  The all-height form
above is convenient: smooth product-cutoff Mellin transforms have
integrable \((1+|t|)\)-weighted absolute value bounded by a fixed
power of \(L\), so there is no boundary-smoothing loss from a sharp bank.

\subsection{The exact target-grid energy}

Let \(S=\{p:p\le y\}\), \(y=\exp(L^{2/5})\), and
\(E_0=\exp(\sqrt L)\), as in Section~\ref{sec:comparison}.
Let \(\xi\) be either centered function in \textup{(C3)}, at a
fixed depth which will be chosen after the required precision.
For a comparison term, define
\[
 F(n)=\psi(n/Y)n^{i\sigma}\xi(n)W_{\mathbf t}^b(n),
 \quad Y=x^{1+o(1)},\quad |\sigma|\le L^C,             \tag{5.22}\label{eq:gr-energy-F}
\]
where the support of \(\psi\) is a fixed compact positive interval.
All its logarithmic derivatives are bounded by fixed powers of \(L\).
The coefficient \(b\) ignores a selected small-bank label, whose count
is exactly one.  This holds for every comparison above.
Reflection changes only fixed character factors and additive signs,
so we state the estimate at a positive endpoint.

\begin{proposition}[Target-grid residue energy]\label{prop:gr-energy}
Suppose that
\[
       q\le xL^B,\qquad
       \frac{Y}{qP_{\max}}\gg_A L^A
                       \quad\text{for every fixed }A,
 \tag{5.23}\label{eq:gr-energy-ratio}
\]
where \(P_{\max}\le\exp(CL^{1/4})\) is the selected small-bank scale.
For every prescribed \(D>0\), a fixed sufficiently deep density model
gives
\[
 E_q(F):=\sum_{a\bmod q}\left|
           \sum_{n\equiv a\ (q)}F(n)\right|^2
                      \ll_D\frac{Y^2}{q}L^{-D}.       \tag{5.24}\label{eq:gr-energy}
\]
The estimate is uniform in all valuations of \(q\).
It also holds with \(F(n)\) multiplied by \(\ee{\beta n}\)
when \(|\beta|Y\) is bounded by a fixed power of \(L\).
\end{proposition}

\begin{proof}
Refine to \(Q=\lcm(q,L_0)\).
Cauchy--Schwarz on the refined classes gives
\(E_q(F)\le(Q/q)E_Q(F)\le L_0E_Q(F)\).
For every gcd profile \(d\mid Q\), put \(r=Q/d\).
Multiplicative orthogonality gives the exact identity
\[
 E_Q(F)=\sum_{d\mid Q}\frac1{\varphi(r)}
       \sum_{\chi\bmod r}\left|\sum_m F(dm)\chi(m)\right|^2.
 \tag{5.25}\label{eq:gr-gcdidentity}
\]
The character zeros impose \((m,r)=1\), so this identity retains
prime powers, including primes dividing \(L_0\). We do not assume
that the endpoint is prime to \(L_0\). Call a character exceptional
when its primitive inducing character has conductor dividing \(L_0\).
There are only a fixed number of such characters for each \(r\).
In particular \(L_0\mid r\) is not required.

We first identify the negligible sets in the norm actually needed.
For a set \(\mathcal B\) and \(|G|\le L^C\), Cauchy--Schwarz in each
residue class gives
\[
 E_Q(G\1_{\mathcal B})
 \le(1+CY/Q)\sum_{n\asymp Y}|G(n)|^2\1_{\mathcal B}(n).
 \tag{5.26}\label{eq:gr-tail-energy}
\]
The factor is \(O(Y/Q)\) here.
The following populations are exponentially small relative to \(Y\),
up to fixed logarithmic factors:
\[
 \begin{array}{ll}
 n_S>E_0:& O(YL^C e^{-cL^{1/10}}),\\
 p\mid n,\ p\mid Q,\ p>y:& O(YL/y),\\
 p\mid n,\ p\mid Q,\ p\text{ in the selected bank}:&
                  O(YL/P_{\min}),\\
 p^2\mid n,\ p\text{ in the selected bank}:&
                  O(Y/P_{\min}).
 \end{array}                                                       \tag{5.27}\label{eq:gr-tailsets}
\]
For the first bound, take \(a=c/\log y\). Counting multiples of
each possible smooth record gives
\[
 \#\{n\le CY:n_S>E_0\}
 \le CY E_0^{-a}\prod_{p\le y}(1-p^{a-1})^{-1}
                 \ll YL^C e^{-cL^{1/10}}.
\]
This estimate uses no property of the indicators. The other bounds
follow by counting multiples,
using \(P_{\min}\ge\exp(cL^{1/8})\).
Thus all these sets are negligible in \eqref{eq:gr-tail-energy}.

For nonexceptional characters remove only selected primes dividing
\(Q\) and selected-prime squares.  Do not yet impose \(n_S\le E_0\),
as that would couple the extracted prime to its cofactor.
The unique selected small-bank prime gives \(n=dpm\).
On a dyadic prime interval \(p\asymp P\), exact dilation invariance
and the independence of \(b\) from this slot make its remaining
coefficient \(\gamma_d(m)\) independent of \(p\), with
\(|\gamma_d(m)|\le L^C\).
It includes the condition that \(m\) has no other selected-bank prime,
and every other bank weight.  Its containing length is
\(M\asymp Y/(dP)\).
Mellin inversion of the product cutoff separates the character sum as
an integral of
\[
 A_\chi(t)B_{d,\chi}(t),\qquad
 A_\chi(t)=\sum_{p\asymp P,\ p\in\mathcal P_g}\chi(p)p^{it},
 \quad
 B_{d,\chi}(t)=\sum_{m\asymp M}\gamma_d(m)\chi(m)m^{it}.
 \tag{5.28}\label{eq:gr-primefactorization}
\]
The omitted Mellin scaling has modulus one.
The transform has polynomial-logarithmic weighted \(L^1\) norm,
including weight \(1+|t|\).
Condition \eqref{eq:gr-energy-ratio} gives
\[
                       M/r\asymp Y/(QP)\gg1,
\]
independently of \(d\).
Elementary character orthogonality therefore yields
\[
 \frac1{\varphi(r)}\sum_{\chi\bmod r}|B_{d,\chi}(t)|^2
 \le(1+CM/r)\sum_{m\asymp M}|\gamma_d(m)|^2
                  \ll \frac{M^2}{r}L^C.             \tag{5.29}\label{eq:gr-cofactorenergy}
\]
Apply Lemma~\ref{lem:gr-primes} to \(A_\chi\) and Minkowski to the
integral in \eqref{eq:gr-primefactorization}.
Its nonexceptional energy at profile \(d\) is
\[
 \ll \frac{Y^2}{Qd}L^C\exp(-cL^{1/8}).               \tag{5.30}\label{eq:gr-nonexceptional}
\]
The prime dyadic decomposition costs only a fixed power of \(L\).
Summing \(d\) is harmless because
\(\sum_{d\mid Q}1/d\le Q/\varphi(Q)\ll\log\log(3Q)\).
There is no number-of-characters or unweighted divisor-count loss.

For exceptional characters use the further restriction
\(n_S=A\le E_0\), and remove \(Q\)-prime divisors above \(y\).
The errors were already controlled in energy.
Now the gcd profile is determined by the smooth record:
\(\gcd(A,Q)=d\), \(n=Av\), \((v,S)=1\).
If \(\chi^*\) is the primitive character inducing an exceptional
\(\chi\bmod r\), then on the retained profile
\[
             \chi(n/d)=\chi^*(A/d)\chi^*(v).
\]
Indeed \(\operatorname{cond}\chi^*\mid r\) and
\((A/d,r)=1\). One must use \(\chi^*\) here: a literal character
modulo \(L_0\) could vanish on a retained bad-prime prefix.
Every other list weight depends only on \(A\), is bounded by \(L^C\),
and the remaining periodic test in \(v\) has fixed period \(L_0\).
The one-variable hypothesis must center this full periodic test;
matching only a scalar principal mean would not suffice.

There is still the condition that \(v\) avoid primes \(p\mid Q\)
above \(y\).  Remove it before applying the fixed-period centering
lemma, extending the same bounded periodic test to those new points.
At each retained record, multiples of such \(p\) number at most
\(CY/(Ap)\) on the containing interval \(0<v\le CY/A\).
Consequently the change in the profile sum is at most
\[
 L^CY\sum_{\substack{A\ S\text{-smooth}\\d\mid A}}\frac1A
                \sum_{\substack{p\mid Q\\p>y}}\frac1p
 \le \frac{Y}{d}L^CV_S^{-1}\frac{\log Q}{y\log y}
 \ll_D (Y/d)L^{-D}.                                  \tag{5.31}\label{eq:gr-avoidance}
\]
We used the exact harmonic identity
\[
 \sum_{\substack{A\ S\text{-smooth}\\d\mid A}}\frac1A
                       =\frac1{dV_S}                 \tag{5.32}\label{eq:gr-record-harmonic}
\]
when \(d\) is \(S\)-smooth; otherwise this restricted profile is empty.

After this extension, hypothesis \textup{(C3)} applies record by
record at length \(Y/A\).
It permits arbitrary bounded fixed-period tests, arbitrary choices
of retained records, and the stated smooth phases.
Choose its logarithmic accuracy after the list and cutoff costs.
Using \eqref{eq:gr-record-harmonic} gives, for every prescribed \(D_0\),
\[
            \left|\sum_mF_{\rm good}(dm)\chi(m)\right|
                         \ll (Y/d)L^{-D_0}.          \tag{5.33}\label{eq:gr-exceptional}
\]
Restoring the avoidance condition costs \eqref{eq:gr-avoidance}.
In \eqref{eq:gr-gcdidentity} the total exceptional contribution is at most
\[
 \frac{Y^2}{Q}L^{-2D_0}
           \sum_{d\mid Q}\frac1d\,\frac{Q/d}{\varphi(Q/d)}
                     \ll \frac{Y^2}{Q}L^{-2D_0+C}.   \tag{5.34}\label{eq:gr-exceptionalenergy}
\]
Choose \(D_0\) large, add the negligible sets, and use the refinement
inequality to prove \eqref{eq:gr-energy}.
Multiplying the cutoff by \(\ee{\beta Yu}\) changes only its fixed
logarithmic derivative bounds when \(|\beta|Y\le L^B\), proving the
last assertion.
\end{proof}

This proof uses only boundedness, exact prime invariance, and the
recordwise periodic means in \textup{(C1)--(C3)}. In particular no
subgroup support, pole order, positivity of an individual character
branch, or unit-support condition at the fixed modulus enters the
energy estimate.

\subsection{From target-grid energy to comparison cancellation}

We now specialize to the reflected family \(k=T=x\) and complete
its comparison estimate with all normalizations visible. The next
section gives the different frequency argument for fixed \(k=h\).
Dyadically localize the two auxiliary products in
\eqref{eq:gr-comparison-term} to \(D_i\asymp H_i\).
Then
\[
 H_i\ge\exp(L^{1/3}),\qquad
 H_i\le\exp(O_m((\log L)L^{2/5})),\qquad
 Y=xH_1H_2.
\]
Insert endpoint dyadic cutoffs before Mellin inversion, as in the
reflection argument.  The resulting ordinary comparison has form
\[
 \frac1Y\int_0^1\widehat F(\theta)\widehat G(-\theta)
                           m(T\theta)\,d\theta,\qquad
 m(\theta)=\E b_1(D_1)b_2(D_2)\ee{\theta D_1D_2}.
 \tag{5.35}\label{eq:gr-fouriercomparison}
\]
External Mellin transforms and their fixed logarithmic costs are
understood here.  The endpoint \(F\) is centered and has exactly the
form in \eqref{eq:gr-energy-F}; \(G\) has squared norm at most
\(YL^C\).  This follows from bounded invariant endpoints and the
pointwise list bound, without any multiplicativity of list coefficients.

By \eqref{eq:gr-productlaw} and
\eqref{eq:gr-integer-bilinear}, for a rational approximation
\(|\theta-a/r|\le r^{-2}\),
\[
 |m(\theta)|^2\ll L^{C_{\rm prod}}
 \left(r^{-1}+H_1^{-1}+\frac{\log(2r)}{H_2}
                    +\frac{r\log(2r)}{H_1H_2}\right). \tag{5.36}\label{eq:gr-multiplier}
\]
Choose an outer cutoff \(B_{\rm out}\), after the finite comparison
family has been fixed.
Dirichlet approximation with denominator ceiling
\(H_1H_2/L^{B_{\rm out}-1}\) makes the multiplier smaller than any
requested inverse logarithmic power unless \(T\theta\) lies within
\(L^{B_{\rm out}}/(H_1H_2)\) of a rational of denominator
\(r\le L^{B_{\rm out}}\).
The corresponding physical frequencies are covered by grids
of denominator \(q=Tr\), with widths
\(O(L^{B_{\rm out}}/(TH_1H_2))=O(L^{B_{\rm out}}/Y)\).
Moreover
\[
 \frac{Y}{qP_{\max}}
 \gg \exp(cL^{1/3}-CL^{1/4})L^{-B_{\rm out}},
\]
so \eqref{eq:gr-energy-ratio} holds uniformly in every gcd profile.

Discrete additive Parseval gives the exact normalization
\[
 \sum_{j\bmod q}|\widehat F(j/q+\beta)|^2
                         =qE_q(F\ee{\beta n}).        \tag{5.37}\label{eq:gr-gridparseval}
\]
Integrating over width \(L^{B_{\rm out}}/Y\) and applying
Proposition~\ref{prop:gr-energy} gives \(YL^{-D}\) for every desired
\(D\), after demanding a larger residue-energy saving.
The polynomial-logarithmic family of \(r\)'s is paid at the same step.
In particular the \(T\) grid points do not create a factor \(T\):
they are already included in \eqref{eq:gr-gridparseval}.
Cauchy--Schwarz and Parseval on the other endpoint bound
\eqref{eq:gr-fouriercomparison} on these grids by an arbitrary inverse
power of \(L\).  Off the grids, use
\eqref{eq:gr-multiplier} and the endpoint norms.
This proves that every comparison in \eqref{eq:gr-comparison} is
negligible to the requested logarithmic accuracy.

Finally \(\Psi(t)=t\Phi(t)\) makes
\[
               \frac{\Psi(|w|/x)}{|w|}
                       =\frac{\Phi(|w|/x)}x.
\]
Thus the harmonic raw correlation is \(1/x\) times the ordinary
banked correlation on the reflected line.
The comparison at scale \(Y\) has exactly the corresponding factor
\(1/Y\).  Paying the finite comparison coefficients therefore proves
the centered banked replacement required in the assembly.

Fixed bad-supported coefficients in \(Aa+Bb=T\) cause no new
analytic condition.  Embed, for example,
\(\widetilde F(n)=\1_{A\mid n}F(n/A)\).
Its exact residue energy is
\[
                  E_q(\widetilde F)=E_{q/(q,A)}(F),
\]
and its core length is divided by the fixed \(A\).
All energy ratios change only by fixed factors.
Bank primes are units at \(A,B\), so dilation invariance and
shared-label removal persist.  This embedding also avoids treating
the physical summands as units modulo \(L_0\) when they are not.
